\newcommand{\ModuleID}{EL-I.2}
\newcommand{\ModuleTitle}{Case, Number, Gender, and Agreement}
\newcommand{\ModuleStage}{Stage I — Foundations}
\newcommand{\ModuleUnit}{I.2}
\newcommand{\ModuleOutcome}{Choose and justify nominal forms from case function, agreement, and government, including syncretic endings}
\newcommand{\ModulePrerequisites}{EL-I.1 mastery: distinguish form from function and produce a basic clause map}
\newcommand{\ModuleExtensionPractice}{Worksheets I.6 and I.7 remain optional remediation or delayed-recall variants}
\newcommand{\ModuleNext}{EL-I.3, First and Second Declension Nouns}
\newcommand{\ModuleMastery}{Parse at least 18 of 20 nominal forms correctly in context, naming both case and function, and repair four of five broken agreement pairs by identifying the controlling noun and the rule that requires each ending.}
\newcommand{\ModuleDelayNote}{}
\newcommand{\ModuleLessonSource}{\CourseRoot/shared/content/volume1/all}
\newcommand{\ModuleExerciseSource}{\CourseRoot/shared/exercises/volume1/all}
\newcommand{\ModuleWorkedBank}{I}
\newcommand{\ModuleExerciseBank}{I}
\newcommand{\ModuleWorksheetVariants}{A,D}
\newcommand{\ModuleScopeNote}{This packet establishes the case-and-agreement decision process; it does not yet teach complete declensional paradigms.}
\newcommand{\ModulePractice}{%
  \SubmoduleExtension
    {The case map}
    {Syncretic endings make a decision process more reliable than memorizing one English cue for each case.}
    {For the project-composed phrase \Latin{servi domini}, list every morphologically possible analysis of each form before context. Then test the two clause frames \Latin{servi dominum laudant} and \Latin{vox servi auditur}.}
    {State the three questions to ask after an ending permits more than one case or number.}
    {Before context, \Latin{servi} can include nominative plural or genitive singular, and \Latin{domini} can include nominative plural or genitive singular. In the first project-composed clause, finite plural \Latin{laudant} supports nominative plural \Latin{servi}, while accusative \Latin{dominum} is object. In the second, singular \Latin{vox} is subject and \Latin{servi} is genitive singular. A complete decision tree asks what the ending permits, what function the clause requires, and what agreement or government confirms it.}
  \SubmoduleExtension
    {Agreement and government}
    {Agreement copies grammatical coordinates; government assigns a required case. Confusing the two hides why an ending changes.}
    {Draw arrows in the lesson's agreement examples from each adjective to its controlling noun and from each governed noun to its governor. Label every arrow either agreement or government.}
    {Without notes, explain why agreeing words need not share the same visible ending.}
    {Every agreement arrow must preserve case, number, and gender; every government arrow must name the governing verb, adjective, preposition, or construction and its required case. A sound retrieval answer notes that nouns and adjectives may belong to different declensional patterns, so grammatical coordinates rather than letter-for-letter endings agree.}
  \SubmoduleExtension
    {Handwritten study and controlled composition}
    {A transformation is useful only when each changed ending can be tied to a changed grammatical coordinate.}
    {Repeat the lesson's recipient transformation, but annotate each changed word with old coordinates, new coordinates, and the rule that forces the change.}
    {Name one form that stayed unchanged and explain why no grammatical coordinate changed.}
    {The annotations must show that plural recipient forms become dative plural and that every agreeing adjective follows its noun's new coordinates. Any unchanged form is acceptable only when its case, number, gender, and function truly remain the same; visual resemblance alone is not a reason.}
}
